\section{Results}
\label{sec:results}

\Cref{tab:main} reports accuracy and cost for every (model, task) pair, \cref{fig:curves} shows test-error
trajectories and behavior\ifcompact.\else, and \cref{fig:tradeoff} places each model on an accuracy-cost plane. \fi The latency sweep is shown in
\cref{fig:latency,fig:ops}.

\begin{table*}[t]
\centering
\caption{Main results (mean\,$\pm$\,std over 5 seeds; held-out test set of 1,000 points). Irreducible test MSE
is 1, 4, and 0.25 for the linear, curved, and chained tasks. Bold = best at shown precision.
$^\dagger$Welch's $t$-test vs.\ Sequential, $p<0.05$. Inference = forward pass over the whole test set
(batch 1,000), single CPU thread.}
\label{tab:main}
\begin{tabular}{@{}lccccc@{}}
\toprule
\textbf{Model} & \textbf{Test MSE} $\downarrow$ & $\boldsymbol{R^2}$ $\uparrow$ & \textbf{Train (ms/step)} $\downarrow$
& \textbf{Inference (ms)} $\downarrow$ & \textbf{Converge (epoch)} $\downarrow$\\
\midrule
\tablebody{tables/tab_main_body.tex}
\bottomrule
\end{tabular}
\end{table*}

\begin{figure*}[t]
\centering
\includegraphics[width=\textwidth]{fig_curves}
\caption{Test MSE per epoch (log scale, mean of five seeds). Hybrids reach low error in far fewer epochs on the
curved and chained tasks. \HB (pink) is fastest to finish and ends lowest on the chained task.}
\label{fig:curves}
\end{figure*}

\subsection{The Control Pair: Identical Math, Different Code}
\label{sec:res-control}
Sequential and Modular compute the same function and, under the same seed, produce the same test MSE on every
task (\cref{tab:main}). Their training step times differ by only \ctrlStepGap\%, confirming that the pilot's
\pilotSpeedup$\times$ ``Modular advantage'' (\cref{sec:pilot}) disappears once the framework and training loop
are held constant. Their batch-1 inference latency, however, differs by \ctrlLatGap\% (\latOneSequential\ vs.\
\latOneModular\us), with non-overlapping 95\% intervals across ten shuffled repeats, so the gap is not run-order
noise. Both dispatch three tensor operations, but Sequential's \texttt{nn.Sequential} container makes
\modSequential\ \texttt{nn.Module} calls per forward pass against Modular's \modModular\ifcompact.\else. The control
pair therefore isolates a second, purely software overhead, telling us how the code organization costs time even when the
math is the same (\cref{sec:res-latency}).\fi

\subsection{Accuracy: Structure Pays Off Only When the Task Has Structure}
\label{sec:res-acc}
\textbf{Linear and curved tasks.} Every model reaches the noise floor on the linear task ($\text{MSE}\approx1$,
$R^2=0.971$), and on the curved task every model finishes within 2\% of the floor of 4. Simple targets
leave nothing for topology to improve upon, so at a 385-parameter budget the arrangement of parameters does not
change the final accuracy. It does change \emph{how fast} the floor is reached. On the curved task \HB converges in
\curvedConvHB\ epochs versus \curvedConvSeq\ for Sequential, about \curvedWallRatio$\times$ less
training time to convergence (\curvedWallHB\,s vs.\ \curvedWallSeq\,s) despite its higher cost per step.

\ifcompact\else  % dropped in the 6-page build
\textbf{Narrow trunks are fragile.} \HC's trunk width also matters for reliability. In a preliminary
configuration with a 17--17 trunk (376 parameters), one of five curved-task seeds settled in a poor solution
(MSE 21.4 vs.\ $\approx$4.0 for the rest). The exact-budget 20--15 trunk used here trains reliably on every
seed, but it still has the largest spread on the curved task.
\fi

\textbf{Chained task.} Here topology matters. \HB, in which the modules mirror the three stages of the target,
reaches \mseChainHB\ test MSE, \chainHBGain\% below Sequential (\mseChainSequential; $p<0.05$). \HA, which has no
knowledge of the stages, still improves by \chainHAGain\% (\mseChainHA; $p<0.05$). Measured against the noise floor
of 0.25, \HB's excess error (\chainExcessHB) is \chainExcessRatio$\times$ smaller than Sequential's
(\chainExcessSeq). Both hybrids also have much smaller variance ($\pm$0.02 vs.\ $\pm$0.12). The
single-paradigm models are still improving at epoch 100 (\cref{fig:curves}c), so part of their gap is slower
convergence rather than a lower ceiling. At a fixed budget, the gap is what a practitioner would get.
Plain breadth does not help: Parallel is the worst model on the chained task (\mseChainParallel).

\ifcompact\else  % dropped in the 6-page build
\begin{figure*}[t]
\centering
\includegraphics[width=\textwidth]{fig_tradeoff}
\caption{Accuracy vs.\ training cost. The $y$-axis is test MSE divided by the best model's MSE on that task (1 =
best). Error bars are $\pm$1 std over the seeds. On simple tasks all models tie on accuracy, so the cheapest
model wins. On the chained task, accuracy can be bought with per-step cost.}
\label{fig:tradeoff}
\end{figure*}
\fi

\subsection{Latency: Operator Count, Not Parameter Count}
\label{sec:res-latency}
\Cref{tab:arch} and \cref{fig:ops} show that, at a nearly constant parameter budget, single inference
latency spans almost 4$\times$, from \latOneModular\us (Modular) to \latOneHB\us (\HB). Parameter count
cannot explain this: all models have 385--387 parameters. Neither can MACs: Parallel and Sequential have the
same 256 MACs, but Parallel is slower, and \HC has the most MACs (350) but is faster than \HA and \HB. What does
explain it is the number of tensor operations dispatched per forward call. A linear fit of batch-1 latency on
operation count gives $R^2=\opsFitRsq$ with a slope of about \opsFitSlope\us per operation. At this scale, each
\texttt{Linear}, \texttt{ReLU}, \texttt{cat}, or add pays a constant framework transmission cost that lessens its
arithmetic. \HB has 13 operations and runs \latOneRatioHB$\times$ slower than Sequential. Counting the ATen
kernels those operations execute gives an equally good fit ($R^2=\kernFitRsq$), so the effect is not an artifact
of how we count.\ifcompact\ Adding \texttt{nn.Module} calls as a second predictor gives $R^2=\opsModFitRsq$ (descriptive only;
six points).\else Adding the number of \texttt{nn.Module} calls as a second predictor explains almost all
remaining variance ($R^2=\opsModFitRsq$; roughly \perOpUs\us per operation and \perModUs\us per module call),
including the Sequential/Modular gap. With six architectures and three fitted coefficients, we treat this
decomposition as being informative rather than predictive.\fi

\begin{figure}[t]
\centering
\includegraphics[width=\columnwidth]{fig_ops_latency}
\caption{Batch-1 inference latency vs.\ tensor operations per forward call. Parameter counts are in
parentheses; error bars are 95\% intervals across ten shuffled repeats. Latency is almost perfectly
explained by the operator count ($R^2=\opsFitRsq$) and is unrelated to the actual parameter count.}
\label{fig:ops}
\end{figure}

\ifcompact
\begin{figure}[t]
\centering
\includegraphics[width=\columnwidth]{fig_latency_batch_1t}
\caption{Inference latency vs.\ batch size (median over 10 shuffled repeats of 1,000 calls, one CPU thread). Curves are flat up to batch
$\approx$64, where overhead dominates and the ranking follows operator count. By batch 4096 the ranking
changes toward the narrow-activation \HC.}
\label{fig:latency}
\end{figure}
\else
\begin{figure*}[t]
\centering
\includegraphics[width=\textwidth]{fig_latency_batch}
\caption{Inference latency vs.\ batch size (median over 10 shuffled repeats of 1,000 calls; bars show the
10th--90th percentile of individual calls). (a) With
one thread, the curves are flat up to batch $\approx$64, where overhead dominates, and the ranking follows
operator count. By batch 4096 it inverts toward the narrow-activation \HC. (b) Four threads rarely help models
this small and sometimes hurt (Parallel at batch 1,000).}
\label{fig:latency}
\end{figure*}
\fi

\textbf{Batch size changes the ranking.} \Cref{fig:latency}\ifcompact\else a\fi\ shows two regimes. Up to a batch of about 64,
latency is nearly flat and ordered by operator count. Past roughly 500 samples, per-element work takes over
and the ordering reshuffles. At batch 4096, \HC becomes the \emph{fastest} model (\latBigHC\us vs.\
\latBigSequential\us for Sequential) despite having the most MACs. This is consistent with memory traffic
having more value than arithmetic: \HC's widest activation is 20 values per sample, while Sequential's is 128 and
Parallel's concatenation writes 128 more. The two models that build intermediate concatenated tensors are the
slowest at batch 4096: \HB (\latBigHB\us) and Parallel (\latBigParallel\us). The latency ranking of these architectures, for this reason, depends on the batch size at of which it
is being measured.

\ifcompact\else  % dropped in the 6-page build (4-thread panel not shown)
\textbf{Threads.} With four occuring threads\ifcompact\else\ (\cref{fig:latency}b)\fi, batch-1 latency is unchanged. At larger batches the
effect is not the same. Sequential improves at batch 4096 (\latBigSequential\ to \threadSeqBigFour\us), but
Parallel at batch 1,000 slows from \threadParallelOne\ to \threadParallelFour\us and \HA at batch 4096 slows to
\threadHABigFour\us. For below kilobyte models, thread synchronization overhead can go over the parallel work.
\fi
